% ============================================================
%  ICRefine — New sections for the full NeurIPS paper
%  Sections: Introduction, Preliminaries, Discussion,
%             Related Work, Conclusion
%
%  These sections slot into the full paper structure:
%    1 Introduction        ← this file
%    2 Preliminaries       ← this file
%    3 ICRefine (Method)   ← findings_draft.tex §1
%    4 Experimental Setup  ← findings_draft.tex §1.6
%    5 Results             ← findings_draft.tex §2.1–2.2
%    6 Analysis            ← findings_draft.tex §2.3–2.7
%    7 Discussion          ← this file
%    8 Related Work        ← this file
%    9 Conclusion          ← this file
%   10 Limitations         ← findings_draft.tex §2.9
%   Appendices             ← findings_draft.tex Appendix
% ============================================================


% ============================================================
\section{Introduction}
\label{sec:intro}
% ============================================================

In-context learning (ICL) achieves strong performance on many NLP tasks by prepending
demonstrations to the inference prompt, yet the quality of those demonstrations depends
heavily on selection, ordering, and format.
Recent work has explored a more compact alternative: replacing a large demonstration bank
with a \emph{cheatsheet}---a concise, structured summary of task heuristics that is
reusable across all test inputs~\citep{csicl2024}.
A single-pass cheatsheet generated from training examples provides a competitive baseline
at a fraction of the per-input token cost of many-shot ICL, and its fixed format makes it
straightforward to deploy across model families at inference time.

A natural next question is whether the quality of a cheatsheet can be improved by
\emph{iterative refinement}: cycling through training failures, patching the rules, and
extending the cheatsheet with structured corrective examples.
Iterative refinement has proven effective in prompt optimization more broadly
\citep{opro2023,ape2022,zhou2022large}, and there is reason to think that a small train
model, if given explicit feedback loops, might produce cheatsheets competitive with those
generated in a single pass by a more capable model.
But a critical question stands apart: does a cheatsheet refined against one model
\emph{transfer} when applied to a different, stronger model at inference time?
Transfer matters because in practice the train model is chosen for efficiency
(cheap API calls during refinement) while the inference model is chosen for quality.
If iterative refinement encodes corrections that are idiosyncratic to the train model's
failure modes, the resulting cheatsheet may be actively harmful to other models.

We introduce \textbf{\textsc{ICRefine}}, a three-phase iterative cheatsheet refinement
pipeline.
\textbf{Phase~0} bootstraps a prior-knowledge (PK) ruleset from 75 training examples in a
single pass.
\textbf{Phase~1} iteratively patches the PK through a sample-and-accept loop that scores
failures, generates candidate revisions, and applies fix-rate and regression gates.
\textbf{Phase~2} partitions remaining training failures by structural error type and
generates one structured case study per coherent failure cluster, using a novel
\texttt{ACTIVATE~IF\,/\,WHY} format that specifies trigger conditions and correction
rationale.
We evaluate \textsc{ICRefine} (train model: gpt-4.1-mini) on five non-ceiling tasks from
BIG-Bench Hard~\citep{bbh}, transferring the resulting cheatsheets to four stronger
evaluation models: GPT-4.1, Claude-3.7-sonnet, Gemini-2.0-flash, and Llama-3.3-70b.
The baseline is the static Cheatsheet-ICL (CS-ICL) approach of \citet{csicl2024},
generated by gpt-4.1 (not mini) in a single pass.

The central finding is a negative one: the full \textsc{ICRefine} pipeline
\emph{underperforms} CS-ICL on all four non-train model families by $-1.0$ to $-3.7$\,pp,
with the deficit statistically distinguishable from zero for GPT-4.1 and Claude.
This failure is not uniform noise: it is \emph{diagnosable}.
We identify three failure modes, characterize each mechanistically, and show that two can
be partially mitigated by targeted interventions---while the task-determined nature of the
failure pattern (not the choice of train model) is the primary obstacle to closing the gap.

\paragraph{Contributions.}
\begin{enumerate}[noitemsep,topsep=2pt]
  \item We propose \textsc{ICRefine}, a three-phase pipeline that combines iterative
        prior-knowledge patching (Phase~1) with partition-based case study generation in a
        structured \texttt{ACTIVATE~IF\,/\,WHY} format (Phase~2).
  \item We present a systematic evaluation of cheatsheet transfer across five model families,
        demonstrating that the transfer gap is \emph{task-determined} rather than
        model-determined: the same failure pattern replicates with both gpt-4.1-mini and
        gemini-2.0-flash as the train model.
  \item We identify and characterize three failure modes of iterative refinement---oracle
        contamination, partition overfitting, and bootstrap convergence limits---and show that
        each is separately diagnosable.
  \item We introduce two targeted interventions: an Evolutionary Algorithm Phase~1 that
        addresses bootstrap convergence ($+6.0$\,pp GS PK accuracy), and oracle-free Phase~2
        generation that reduces oracle contamination ($+1.9$--$+3.8$\,pp CJ accuracy across
        all models).
        The combined EA+no-oracle pipeline remains below CS-ICL, providing a bounded negative
        result useful to practitioners considering iterative cheatsheet refinement.
\end{enumerate}


% ============================================================
\section{Preliminaries}
\label{sec:prelim}
% ============================================================

\paragraph{Task setting.}
We consider multi-class classification tasks where each input $x$ is mapped to a label
$y \in \mathcal{Y}$ via an LLM prompted with a fixed instruction.
We have access to a labeled training set $\mathcal{D} = \{(x_i, y_i)\}_{i=1}^{N}$ and
evaluate on a held-out test set.
At inference time, no labeled demonstrations are provided in the prompt; the model
generates a chain-of-thought and produces a final verdict.
All evaluations use \emph{reasoning-first} (RF) scoring: the model writes its reasoning
before emitting \texttt{VERDICT:~(X)}.

\paragraph{Cheatsheet.}
A \emph{cheatsheet} $\mathcal{C} = (\text{PK}, \mathcal{S})$ consists of two components
prepended to every test prompt:
(i) a \textbf{prior knowledge} (PK) block---a free-form natural-language ruleset
($\leq$12\,000 characters) encoding general task heuristics and common error corrections;
and (ii) a set of \textbf{case studies} $\mathcal{S}$, each in \texttt{ACTIVATE~IF\,/\,WHY}
format: a set of structural trigger predicates over the input, a correction rationale, and
annotated support examples.
A static cheatsheet ($|\mathcal{S}| = 0$) is the CS-ICL baseline of \citet{csicl2024};
\textsc{ICRefine} produces cheatsheets with both components.

\paragraph{Train and eval models.}
We distinguish the \textbf{train model} $M_\text{tr}$, used for all pipeline roles during
refinement (scoring, rule-patching, case study generation), from \textbf{evaluation models}
$\{M_\text{ev}\}$, which consume the cheatsheet at inference time.
We say a cheatsheet \emph{transfers} if it improves $M_\text{ev}$ accuracy above the CS-ICL
baseline when $M_\text{ev} \neq M_\text{tr}$.
Transfer is non-trivial because corrections derived from $M_\text{tr}$'s failure modes may
not match $M_\text{ev}$'s failure modes, and may actively mislead a more capable model.


% ============================================================
\section{Discussion}
\label{sec:discussion}
% ============================================================

Our results point toward several generalizable principles for iterative cheatsheet design
that extend beyond the specific pipeline and task suite studied here.

\paragraph{Principle 1: Task failure structure, not model identity, governs transferability.}
The transfer deficit of \textsc{ICRefine} over CS-ICL is nearly identical whether the
train model is gpt-4.1-mini or gemini-2.0-flash ($-3.4$ vs.\ $-3.7$\,pp for GPT-4.1;
$-3.7$ vs.\ $-2.9$\,pp for Claude).
The locus of harm is also consistent across training directions: CJ case studies are
universally harmful, GS case studies benefit models that share the train model's SVG parsing
failures.
This suggests that the key design variable is \emph{failure mode coverage}---whether the
failure modes encoded in a case study are broadly shared across model families---rather than
the identity of the model used to generate the cheatsheet.
Case study transferability is therefore a \emph{property of the task} (specifically, whether
its failures are structurally universal or model-specific) rather than a property of the
training regime.
Practitioners should audit whether target-model failure modes overlap with train-model failure
modes before deploying iteratively refined cheatsheets.

\paragraph{Principle 2: Oracle context helps where failure modes are principled, hurts where they are idiosyncratic.}
Oracle injection (gold chain-of-thought alongside each training failure) was intended to help
the train model identify better correction rules.
It produces opposite effects on our two most informative tasks: GS case studies generated
with oracle access encode structurally transferable geometric rules (vertex counting, arc
classification), while CJ case studies encode idiosyncratic reasoning traces for one specific
causal scenario.
The oracle is not itself to blame---the CJ E3 intervention shows that oracle-free generation
narrows but does not close the gap, because heterogeneous failure distribution is a
co-cause.
The design implication is that oracle context should be reserved for tasks where the failure
mode admits a principled, item-independent correction rule; for tasks with dispersed,
heterogeneous failure distributions, oracle access risks memorising rather than abstracting.

\paragraph{Principle 3: A strong bootstrap with a capable generator can dominate iterative refinement.}
CS-ICL, generated by gpt-4.1 in a single pass over all 100 training examples, consistently
outperforms \textsc{ICRefine} (train model: gpt-4.1-mini) on the four non-train evaluation
models.
Part of this gap reflects generator quality: gpt-4.1 (not mini) produces the CS-ICL
cheatsheet, so a mini-generated CS-ICL would remove this confound.
But a deeper point holds independently of generator quality: a model reading all 100
training examples in a single, unconstrained pass may extract more broadly applicable
heuristics than an iterative pipeline that focuses each revision on the current failure set.
Iterative refinement is well-suited to structured, enumerable failure modes (systematic
logical inference errors, specific SVG parsing bugs), but a capable model in a single pass
can capture holistic task intuitions that no sequence of local patches recovers.
Iterative refinement should therefore be evaluated not just against its own prior state but
against a strong single-pass generator as a baseline.

\paragraph{Principle 4: Partition overfitting is a structural failure of greedy partition search.}
Phase~2's greedy partition search re-encounters the same failure type across multiple bins
and generates redundant case studies without detecting the overlap (formal fallacies: three
of four case studies address illicit conditional conversion; CJ seeds~2--3 converge on
joint-AND causation from three different angles).
This is not a failure of the generation model---it is a failure of the search procedure.
Deduplication via ACTIVATE~IF similarity checks before acceptance, and diversity-promoting
objectives in case study selection, would mechanically address this issue.
The parallel with best-of-$N$ selection failures in the size ablation (test ordering
reverses from train on GS) suggests a broader lesson: train-time acceptance metrics
(fix-rate, regression-rate) are necessary but not sufficient conditions for test-time
generalization of case studies.

\paragraph{Principle 5: Evolutionary search adds value precisely when greedy search stalls.}
EA Phase~1 produces meaningful improvements on the two tasks where sequential patching
stalls (GS: $+6.0$\,pp PK accuracy; CJ: $+5.4$\,pp PK accuracy) and is within noise on
tasks that sequential patching already handles (Snarks, DQ).
Critically, EA's improvement on GS \emph{eliminates} the need for Phase~2 case studies
entirely: the acceptance gate rejects all candidates because the EA prior knowledge is
already comprehensive enough.
This suggests a practical heuristic---use single-candidate patching until stall, then switch
to population-based search---and motivates a cascade design where the pipeline monitors
patch acceptance rates and escalates search strategy only when needed.
Population-based search is more expensive, but the additional cost is only incurred where it
is causally necessary.


% ============================================================
\section{Related Work}
\label{sec:related}
% ============================================================

\paragraph{In-context learning.}
In-context learning (ICL) emerged as a central capability of large language models with
GPT-3~\citep{brown2020gpt3}, which showed that a few labeled examples prepended to a prompt
suffice for strong zero-shot-equivalent task performance.
Subsequent work established that demonstration selection, ordering, and format can each
produce large accuracy swings independent of model capability~\citep{zhao2021calibrate,lu2022order}.
Chain-of-thought (CoT) prompting~\citep{wei2022cot} extended ICL to multi-step reasoning
by eliciting intermediate reasoning traces, and many-shot ICL~\citep{agarwal2024manyshot}
showed that scaling the number of demonstrations continues to improve performance up to the
context window.
\textsc{ICRefine} operates in the complementary regime: rather than maximizing the number of
demonstrations, it compresses task knowledge into a fixed-length cheatsheet that is prepended
to all test inputs, reducing per-input token cost while retaining task-specific guidance.

\paragraph{Cheatsheet and prompt distillation.}
\citet{csicl2024} introduced Cheatsheet-ICL (CS-ICL), the direct baseline for our work:
they show that a single-pass natural-language ruleset generated from training demonstrations
matches or exceeds many-shot ICL accuracy at a fraction of the inference cost.
Related work on prompt distillation~\citep{li2023promptdistillation} and task-specific
instruction generation~\citep{ape2022,zhou2022large} similarly attempts to compress
demonstration knowledge into reusable prompts.
\textsc{ICRefine} extends CS-ICL with iterative refinement and a structured case study
component, and contributes a cross-model transfer evaluation that CS-ICL does not perform.

\paragraph{Automatic prompt optimization.}
A growing body of work treats prompt construction as an optimization problem.
Automatic Prompt Engineer (APE)~\citep{ape2022} generates and selects candidate
instructions using LLM proposals and a scoring function.
OPRO~\citep{opro2023} frames prompt optimization as black-box optimization, using
past (prompt, score) pairs as meta-prompts to guide refinement.
ProTeGi~\citep{pryzant2023protegi} refines prompts by generating natural-language ``gradient''
critiques of failures and applying them as directed edits.
RLPrompt~\citep{deng2022rlprompt} uses reinforcement learning to search the discrete prompt
space.
\textsc{ICRefine}'s Phase~1 is most closely related to ProTeGi-style failure-directed
editing, but differs in applying gated accept/reject logic rather than gradient-style
accumulation, and in operating on a structured cheatsheet (PK + case studies) rather than a
free-form instruction.
Our EA Phase~1 variant adds a population-based search strategy that is not present in any of
these single-candidate approaches.

\paragraph{Iterative and self-improving NLP systems.}
Iterative self-refinement~\citep{madaan2023selfrefine} uses the model's own feedback to
revise its output across multiple rounds, while Reflexion~\citep{shinn2023reflexion}
accumulates verbal experience from failures in episodic memory.
These approaches refine individual \emph{outputs} at inference time; \textsc{ICRefine}
instead refines a \emph{shared cheatsheet} at train time that is then applied
read-only at inference time.
Constitutional AI~\citep{bai2022constitutional} and RLHF-adjacent iterative approaches
~\citep{ouyang2022instructgpt} similarly use feedback to improve behavior, but target model
weights rather than inference-time prompts.
The sample-and-accept gating in Phase~1 is related to the self-play and candidate filtering
used in STaR~\citep{zelikman2022star}, which iteratively bootstraps reasoning chains from
correct self-generated rationales.

\paragraph{Cross-model prompt transfer.}
Several papers have noted that prompts optimized for one model often transfer imperfectly to
others~\citep{lester2021power,logan2021cutting}.
Soft-prompt tuning~\citep{lester2021power} encodes task knowledge in continuous token
embeddings, which are model-specific by construction; hard (discrete) prompt optimization
\citep{ape2022,opro2023} is more portable in principle.
Recent work on model-agnostic prompting~\citep{chen2023modelagnostic} studies conditions
under which task instructions generalize across model families.
Our contribution to this thread is a controlled cross-model transfer evaluation showing that
cheatsheet quality---measured on the train model---is a poor predictor of transfer quality,
and that the failure mode structure of the \emph{task} (not the model) is the primary
determinant of transferability.


% ============================================================
\section{Conclusion}
\label{sec:conclusion}
% ============================================================

We introduced \textsc{ICRefine}, a three-phase iterative pipeline for cheatsheet
construction, and evaluated its ability to transfer task-specific knowledge from a small
train model to stronger evaluation models on BIG-Bench Hard.
The central finding is a bounded negative result: despite producing cheatsheets with
substantially higher training accuracy than the static CS-ICL baseline, the full pipeline
underperforms CS-ICL on all four non-train evaluation model families, with statistically
significant deficits for GPT-4.1 and Claude.
We characterize the failure as task-determined rather than model-determined: three
mechanistically distinct failure modes---oracle contamination, partition overfitting, and
bootstrap convergence limits---each map to specific tasks and persist regardless of which
model generates the cheatsheet.

Two targeted interventions---evolutionary Phase~1 and oracle-free Phase~2---partially
address two of the three failure modes.
EA Phase~1 closes the bootstrap convergence gap on geometric\_shapes ($+6.0$\,pp PK
accuracy) and eliminates the need for Phase~2 case studies on that task entirely.
Oracle-free Phase~2 reduces but does not eliminate the CJ transfer gap ($+1.9$--$+3.8$\,pp),
with the residual attributable to heterogeneous failure distribution that is orthogonal to
oracle access.
The combined EA+no-oracle pipeline remains below CS-ICL for GPT-4.1 and Llama, suggesting
that the bottleneck is the difficulty of case study generalization under diverse failure
distributions, not any single addressable design flaw.

These results reframe iterative cheatsheet refinement as a problem whose difficulty is
determined by task structure: tasks with concentrated, structurally universal failure modes
are amenable to refinement and case study generation; tasks with dispersed, model-specific,
or idiosyncratic failures are not.
A practical upshot for practitioners is that a strong single-pass generator applied to all
training examples should serve as the primary baseline before investing in iterative
refinement, and that train-model accuracy gains during refinement do not reliably predict
cross-model transfer quality.
We release our pipeline code, cheatsheets, and evaluation logs to support follow-on work.
